\chapter{Evaluation Model} %Abstract
Expressions in EDADSL are evaluated according to standard mathematical precedence rules. The evaluation model is eager (strict), meaning all sub-expressions are evaluated before the operator is applied.

\section{Operator Precedence}
% Reference to Chapter 7 (Expressions)
The complete operator precedence is defined in Section~\ref{sec:operators}.

\section{Type Coercion}
Type coercion follows the hierarchy: \texttt{int} $<$ \texttt{real} $<$ \texttt{complex}.
\begin{itemize}
\item \texttt{int} coerces to \texttt{real}: \texttt{5} becomes \texttt{5.0}
\item \texttt{real} coerces to \texttt{complex}: \texttt{3.14} becomes \texttt{3.14 + 0i}
\item \texttt{int} coerces to \texttt{complex}: \texttt{5} becomes \texttt{5 + 0i}
\item \texttt{bool} coerces to \texttt{int}: \texttt{True} becomes \texttt{1}
\item \texttt{bool} coerces to \texttt{real}: \texttt{True} becomes \texttt{1.0}
\item \texttt{bool} coerces to \texttt{complex}: \texttt{True} becomes \texttt{1 + 0i}
\item Mixed numeric operations promote to the wider type: \texttt{int + real} $\rightarrow$ \texttt{real}, \texttt{real + complex} $\rightarrow$ \texttt{complex}
\item Coercion is never lossy in the upward direction
\item Downward coercion (\texttt{complex} $\rightarrow$ \texttt{real}, \texttt{real} $\rightarrow$ \texttt{int}) requires explicit casting and may raise errors
\end{itemize}


\chapter{Execution Order}  %Concrete

\section{Sequential Execution}
Program statements are executed sequentially from the entry point (\texttt{start}).

\section{Declaration Order}
% TODO: Define whether declarations are processed before execution
% Uncertain about:
% - Can parts be referenced before they are declared?
% - How does the order of elec.part and elec.connect statements affect the result?


\chapter{Constraint Resolution}
Physical constraints are resolved by the solver according to the priority system defined in Section~\ref{sec:physical-constraints}.

\section{Conflict Resolution}
When constraints conflict, the solver applies the higher-priority constraint and optimizes the lower-priority one. See Section~\ref{sec:solver-model} for the complete resolution rules.


\chapter{Type System Rules}
EDADSL uses a dynamic type system with distinct types for electrical and non-electrical objects.

\section{Type Categories}
\begin{itemize}
\item Generic types: Boolean, Integer, Real, Complex, String, List, Set
\item Electrical types: Part, Pin, Connection, Net
\end{itemize}

\section{Type Safety}
% TODO: Define formal type safety guarantees
% Uncertain about:
% - Whether sets can contain mixed types
% - How function argument type checking works
% - Runtime vs compile-time type checking




\part{Core Formal Model}
